% !TeX program = xelatex
% Bien dich tu reports/: xelatex -interaction=nonstopmode -halt-on-error main.tex
\documentclass[12pt,a4paper]{article}
\usepackage{fontspec}
\IfFontExistsTF{Noto Serif}{\setmainfont{Noto Serif}}{\setmainfont{DejaVu Serif}}
\IfFontExistsTF{Noto Sans}{\setsansfont{Noto Sans}}{\setsansfont{DejaVu Sans}}
\IfFontExistsTF{DejaVu Sans Mono}{\setmonofont{DejaVu Sans Mono}}{\setmonofont{Latin Modern Mono}}
\usepackage[margin=2.5cm]{geometry}
\usepackage{amsmath,graphicx,array,longtable,booktabs,tabularx}
\usepackage{xcolor,fvextra}
\usepackage{tikz}
\usetikzlibrary{arrows.meta,positioning}
\usepackage[unicode,hidelinks]{hyperref}
\usepackage{caption}
\renewcommand{\contentsname}{Mục lục}
\renewcommand{\figurename}{Hình}
\renewcommand{\tablename}{Bảng}
\renewcommand{\appendixname}{Phụ lục}
\setlength{\parindent}{1.2em}
\setlength{\parskip}{0.35em}
\setlength{\emergencystretch}{3em}
\newcommand{\code}[1]{\texttt{\detokenize{#1}}}
\newcommand{\src}[2]{\par\noindent{\small\textbf{Đối chiếu:} \path{#1}; \code{#2}.}\par}
\newcommand{\todo}[1]{\par\noindent\textcolor{red!65!black}{\textbf{TODO:} #1}\par}
\newcommand{\reportfigure}[4]{%
  \begin{figure}[htbp]\centering
  \IfFileExists{#1}{\includegraphics[width=.92\linewidth,height=.42\textheight,keepaspectratio]{#1}}{%
    \fbox{\parbox[c][4.0cm][c]{.87\linewidth}{\centering\textbf{Ảnh cần bổ sung}\\[.6em]#4\\[.6em]\small\path{#1}}}}
  \caption{#2}\label{#3}\end{figure}}
\DefineVerbatimEnvironment{Code}{Verbatim}{fontsize=\footnotesize,breaklines=true,breakanywhere=true,frame=single}
\begin{document}
\begin{titlepage}
\centering
\vspace*{2cm}
{\Large BÁO CÁO THỰC HÀNH\par}
\vspace{1cm}
{\LARGE\bfseries Bài thực hành 03\\[.5em]LLM Skill Planning với Gripper và Camera\par}
\vspace{2cm}
{\large Sinh viên: Lê Hồng Quang\\[.5em]MSSV: 23020757\par}
\vspace{1.5cm}
{\large Môi trường: ROS 2, MoveIt 2 và Gazebo\par}
\vfill
{\small Báo cáo đối chiếu mã nguồn tại \path{~/ur_ws/src/ur3_vision_planning}.\\
Ngày tổng hợp: 05/10/2026.\\
Các kết quả chưa có bằng chứng được đánh dấu TODO.\par}
\end{titlepage}
\tableofcontents\newpage

\section{Giới thiệu}
Bài thực hành xây dựng hệ thống điều khiển thao tác vật bằng câu lệnh ngôn ngữ tự nhiên trong Gazebo. Môi trường gồm robot UR3, gripper song song hai ngón, camera RGB-D đặt phía trên, bàn thao tác, năm block màu và ba vùng đặt vật. Người dùng chỉ định một vật và một vùng đích; hệ thống nhận diện trạng thái hiện tại, tạo chuỗi thao tác và kiểm tra trước khi điều khiển robot.

LLM được sử dụng ở cấp ý định, không sinh góc khớp hoặc quỹ đạo. Trong luồng Bài 03, kết quả LLM chỉ gồm tên vật và zone đích. Logic xác định trong chương trình xử lý occupancy, chọn vị trí tạm, mở rộng kế hoạch và kiểm tra thứ tự skill. Khi vùng đích có vật khác, kế hoạch phải dời vật đang chiếm ra trước. Bài 03 không có nhiệm vụ cá nhân hóa theo mã sinh viên; thông tin sinh viên chỉ dùng tại trang bìa.

Báo cáo phân biệt ba mức: \textbf{xác nhận từ mã nguồn} nghĩa là có triển khai; \textbf{bằng chứng chạy thử} phải gắn với log hoặc ghi nhận cụ thể; \textbf{cần kiểm tra thêm} là phần chưa đủ dữ liệu thực nghiệm. Lần tổng hợp này chỉ đọc nguồn và tài liệu, không chạy simulator, chuyển động robot, gọi API LLM hoặc đọc tệp cấu hình bí mật. Việc mô tả lệnh chạy trong phụ lục không có nghĩa các lệnh đã được thực hiện trong lần phân tích này.

\section{Công nghệ sử dụng}
\begin{tabularx}{\linewidth}{@{}p{.27\linewidth}X@{}}\toprule
Công nghệ & Vai trò trong chương trình\\\midrule
ROS 2 Jazzy, Python, \code{rclpy} & Node, parameter, topic, service và action; workspace có \path{/opt/ros/jazzy}.\\
Gazebo, \code{ros_gz_sim}, \code{ros_gz_bridge} & Mô phỏng vật lý, spawn vật, camera RGB-D, chuyển ảnh và đồng hồ sang ROS. README đề cập Gazebo 8.\\
MoveIt 2, \code{ur_moveit_config} & IK, FK, Cartesian path, lập kế hoạch khớp, collision scene và thực thi quỹ đạo.\\
\code{gz_ros2_control}, UR controllers & Điều khiển tay máy và gripper trong mô phỏng.\\
OpenCV, NumPy, \code{cv_bridge}, TF2 & Phân đoạn HSV, xử lý depth, lọc tọa độ và chuyển hệ quy chiếu.\\
9Router, HTTP tương thích OpenAI & Chuyển câu lệnh tiếng Việt/Anh thành JSON mục tiêu.\\
Xacro, YAML, pytest & Mô tả môi trường, cấu hình và kiểm thử logic bằng fixture.\\\bottomrule
\end{tabularx}
Các phụ thuộc trực tiếp được khai báo trong \path{src/ur3_vision_planning/package.xml}; phụ thuộc mô phỏng UR được khai báo trong \path{src/ur_simulation_gz/ur_simulation_gz/package.xml}. Hai manifest cần được đọc cùng với import và launch, vì controller, MoveIt action server và cấu hình UR còn đến từ các package được include.

\section{Kiến trúc hệ thống}
\begin{figure}[htbp]\centering
\begin{tikzpicture}[node distance=5mm, every node/.style={font=\small},
box/.style={draw,rounded corners,align=center,text width=10.8cm,minimum height=8mm},
arr/.style={-{Latex},thick}]
\node[box] (u) {Câu lệnh người dùng};
\node[box,below=of u] (l) {9Router / StructuredGoalPlanner: JSON \{object, target\_zone\}};
\node[box,below=of l] (e) {Environment Manager: snapshot camera mới, occupancy, vị trí tạm};
\node[box,below=of e] (v) {Expanded plan $\rightarrow$ DryRunPlanValidator};
\node[box,below=of v] (p) {VisionRobotSkills: TF, IK, Cartesian grasp/lift, HOME planning-only};
\node[box,below=of p] (x) {execute\_validated\_plan: pick, place, camera verification, home};
\node[box,below=of x] (m) {MoveIt 2 $\rightarrow$ trajectory guard $\rightarrow$ ExecuteTrajectory};
\node[box,below=of m] (g) {UR3 + gripper + DetachableJoint trong Gazebo};
\foreach \a/\b in {u/l,l/e,e/v,v/p,p/x,x/m,m/g}\draw[arr] (\a)--(\b);
\end{tikzpicture}
\caption{Luồng xử lý chính. Camera cung cấp trạng thái trước lập kế hoạch và xác nhận sau đặt vật.}\label{fig:architecture}
\end{figure}

Trong \code{resolve_command}, yêu cầu HTTP được thực hiện trên một thread riêng để node tiếp tục nhận camera. Sau khi JSON mục tiêu hợp lệ, \code{fresh_snapshot} yêu cầu bản tin môi trường mới, thay vì dùng snapshot trước khi nhập lệnh hoặc trước khi gọi LLM. \code{EnvironmentPlanner.resolve_goal} mở rộng mục tiêu thành kế hoạch; \code{validate_plan} kiểm tra lại mục tiêu và vị trí tạm đã chọn.

LLM quyết định \emph{người dùng muốn di chuyển vật nào đến zone nào}. Việc chọn occupant cần dời, vị trí tạm, chuỗi pick/place/home và tọa độ thực thi đều do chương trình xử lý. MoveIt tạo chuyển động từ các mục tiêu hình học, trạng thái khớp và collision scene. Camera xác nhận kết quả place trước bước tiếp theo. Không có topic ROS dành riêng cho việc truyền goal LLM sang executor trong luồng này: các module trao đổi đối tượng Python và gọi hàm trong cùng tiến trình.
\src{src/ur3_vision_planning/ur3_vision_planning/natural_language_planner.py}{StructuredGoalPlanner, resolve_command, NaturalLanguagePlanner.fresh_snapshot}
\src{src/ur3_vision_planning/ur3_vision_planning/plan_execution.py}{execute_validated_plan}

\section{Cấu trúc chương trình và các node/package chính}
\subsection{Tổ chức thư mục và phần kế thừa}
Package \code{ur3_vision_planning} là package \code{ament_python}. \code{setup.py} cài module và đăng ký console scripts; \code{setup.cfg} quy định vị trí executable. Cấu trúc chính:
\begin{Code}
src/ur3_vision_planning/
  launch/llm_robot.launch.py
  config/{scene,camera,perception,temporary_positions,execution}.yaml
  config/ur3_gripper_controllers.yaml
  config/student_config.yaml        # cấu hình legacy; không phải yêu cầu bài
  worlds/vision_world.sdf.xacro
  urdf/ur3_table_mount.urdf.xacro
  urdf/simple_parallel_gripper.urdf.xacro
  urdf/robotiq_2f_85_macro_no_mimic.urdf.xacro
  ur3_vision_planning/              # các module Python
  test/                            # 12 tệp kiểm thử
  README.md, setup.py, setup.cfg, package.xml
src/ur3_llm_control/                # package Bài 02 độc lập
src/ur_simulation_gz/ur_simulation_gz/
  launch/ur_sim_control.launch.py
  launch/ur_sim_moveit.launch.py
src/ur3_circle/                    # không thuộc pipeline chính của Bài 03
\end{Code}

Đối chiếu hai package cho thấy Bài 03 giữ nền tảng HTTP planner, validator legacy, scene publisher, skill pick/place/home và mô tả robot/gripper từ Bài 02. Phần bổ sung gồm world RGB-D, hai block xanh lá và tím, perception camera, Environment Manager, planner mục tiêu ngôn ngữ tự nhiên, executor có camera verification, precheck và kiểm tra liên tục nhánh nghiệm/quỹ đạo. Mô hình đang được include là \code{simple_parallel_gripper}; sự hiện diện của tệp macro Robotiq không chứng minh gripper Robotiq đang được dùng.
\src{src/ur3_llm_control/setup.py}{console_scripts}
\src{src/ur3_vision_planning/setup.py}{console_scripts, data_files}
\src{src/ur3_vision_planning/urdf/ur3_table_mount.urdf.xacro}{xacro:ur_robot, xacro:simple_parallel_gripper}

\subsection{Executable và ROS node thực tế}
\begin{longtable}{@{}p{.29\linewidth}p{.25\linewidth}p{.39\linewidth}@{}}
\toprule Executable & Node/class & Chức năng thực tế\\\midrule\endhead
\code{camera_perception} & \code{camera_perception}; \code{CameraPerception} & Nhận RGB-D, xuất debug và JSON trạng thái.\\
\code{scene_publisher} & \code{scene_publisher}; \code{ScenePublisher} & Khởi tạo collision table/cubes và ACM trong MoveIt.\\
\code{environment_manager_test} & \code{environment_manager_test}; \code{EnvironmentManager} & Nhận camera, resolve goal từ parameter, in plan và thoát; không chuyển động.\\
\code{natural_language_planner} & \code{natural_language_planner} và \code{llm_planner} & Node manager cộng node HTTP; dry-run mặc định, có nhánh precheck/execution.\\
\code{motion_precheck} & \code{NaturalLanguagePlanner} & Entrypoint nhận object/zone từ parameter, bỏ qua LLM, bắt buộc precheck-only. Không phải class node tên MotionPrecheck.\\
\code{llm_planner} & \code{llm_planner}; \code{LLMPlanner} & Planner legacy trả danh sách skill, giới hạn ba màu.\\
\code{skill_executor} & \code{llm_planner}, \code{robot_skills} & Entrypoint interactive legacy; dùng \code{TaskValidator} và \code{VisionRobotSkills}, không phải pipeline camera expanded plan.\\
\bottomrule
\end{longtable}
\code{RobotSkills} kế thừa \code{Node}; \code{VisionRobotSkills} kế thừa nó và giữ tên node \code{robot_skills}, chỉ được tạo ở nhánh motion của pipeline mới. \code{EnvironmentPlanner}, \code{DryRunPlanValidator}, \code{TaskValidator}, \code{CameraPlacementConfirmation} là class hỗ trợ, không phải ROS node. \code{ContinuousMotionMixin}, \code{home_precheck.py}, \code{joint_continuity.py}, \code{scene_config.load_scene} cũng không đăng ký executable riêng.

Launch chính include mô phỏng UR và MoveIt; tạo perception, scene publisher, gripper spawner, hai static TF và bridge camera. Các node hạ tầng còn có \code{robot_state_publisher}, \code{controller_manager/spawner}, \code{ros_gz_sim/create}, bridge \code{/clock}, RViz và MoveIt từ launch phụ thuộc. Tên node cụ thể của MoveIt/runtime cần đối chiếu bằng \code{ros2 node list} khi chạy; không đồng nhất tên package, executable và tên node.
\src{src/ur3_vision_planning/launch/llm_robot.launch.py}{launch_setup, generate_launch_description}
\src{src/ur_simulation_gz/ur_simulation_gz/launch/ur_sim_control.launch.py}{launch_setup}
\src{src/ur3_vision_planning/ur3_vision_planning/motion_precheck.py}{main, repeat_summary}

\subsection{Giao tiếp ROS 2}
\begin{longtable}{@{}p{.45\linewidth}p{.48\linewidth}@{}}\toprule
Topic & Kiểu và mục đích\\\midrule\endhead
\code{/camera/image_raw} & \code{sensor_msgs/msg/Image}, RGB từ Gazebo.\\
\code{/camera/depth/image_raw} & \code{sensor_msgs/msg/Image}, depth đã căn chỉnh.\\
\code{/camera/camera_info} & \code{sensor_msgs/msg/CameraInfo}, ma trận nội tại và frame.\\
\code{/vision/debug_image} & \code{sensor_msgs/msg/Image}, bbox, tên, tọa độ và zone.\\
\code{/vision/environment_state} & \code{std_msgs/msg/String}, JSON môi trường; manager subscribe.\\
\code{/joint_states} & \code{sensor_msgs/msg/JointState}, khớp tay/ngón.\\
\code{/simple_gripper_controller/commands} & \code{std_msgs/msg/Float64MultiArray}, hai vận tốc ngón.\\
\code{/planning_scene} & \code{moveit_msgs/msg/PlanningScene}; publisher được tạo, đường khởi tạo scene chủ yếu qua service.\\
\code{/robot_description}, \code{/robot_description_semantic} & \code{std_msgs/msg/String}, fallback mô hình runtime cho HOME.\\
\code{/tf}, \code{/tf_static}, \code{/clock} & \code{tf2_msgs/msg/TFMessage}, \code{rosgraph_msgs/msg/Clock}.\\\bottomrule
\end{longtable}
Service MoveIt chính gồm \code{/compute_ik} (\code{GetPositionIK}), \code{/compute_fk} (\code{GetPositionFK}), \code{/compute_cartesian_path} (\code{GetCartesianPath}), \code{/check_state_validity} (\code{GetStateValidity}), \code{/get_planning_scene} (\code{GetPlanningScene}), \code{/apply_planning_scene} (\code{ApplyPlanningScene}) và \code{/plan_kinematic_path} (\code{GetMotionPlan}), đều thuộc \code{moveit_msgs/srv}. HOME còn dùng \code{/move_group/list_parameters}, \code{/move_group/get_parameter_types}, \code{/move_group/get_parameters}, với các kiểu tương ứng trong \code{rcl_interfaces/srv}.

Action client nền gồm \code{/move_action} (\code{moveit_msgs/action/MoveGroup}) và \code{/execute_trajectory} (\code{moveit_msgs/action/ExecuteTrajectory}). Adapter Bài 03 thực hiện plan-only và gửi quỹ đạo đã kiểm tra qua \code{ExecuteTrajectory}. Controller tay cấu hình \code{scaled_joint_trajectory_controller}; action \code{FollowJointTrajectory} thuộc hạ tầng controller, không phải đầu ra của LLM.

Các topic \code{/ur3_vision_planning/grasp/<object>/attach}, \code{detach}, \code{state} và service \code{/world/empty/entity/system/add} là \emph{Gazebo Transport}, không phải ROS 2 topic/service được bridge trong launch. Lệnh attach/detach phát \code{gz.msgs.Empty}; discovery Gazebo 8 phân biệt attach subscriber generic \code{google.protobuf.Message} với detach subscriber \code{gz.msgs.Empty}; state được đọc và phân tích chuỗi \code{attached}/\code{detached}.
\src{src/ur3_vision_planning/ur3_vision_planning/robot_skills.py}{RobotSkills.__init__, set_gazebo_attachment, attachment_subscriber_present}
\src{src/ur3_vision_planning/ur3_vision_planning/vision_robot_skills.py}{VisionRobotSkills.__init__, _home_model}

\section{Camera và nhận diện trạng thái môi trường}
\subsection{Nguồn dữ liệu và nhận diện năm block}
World \code{empty} chứa camera \code{overhead_camera}, sensor \code{overhead_rgbd} loại \code{rgbd_camera}, render engine \code{ogre2}. Cấu hình đặt camera tại $(0.38,0,1.15)$ m, ảnh $640\times480$, tốc độ cấu hình 15 Hz, góc nhìn ngang $60^\circ$. Launch bridge các Gazebo topic \code{/camera/image}, \code{/camera/camera_info}, \code{/camera/depth_image} sang ba ROS topic camera. Static TF nối \code{world}, \code{camera_link}, \code{camera_optical_frame}.

\code{CameraPerception.process_pending} ghép RGB với depth gần timestamp nhất trong queue dài 8; chênh lệch tối đa 0.08 s. \code{process_pair} yêu cầu frame và kích thước RGB/depth/CameraInfo trùng nhau, tiêu cự dương và distortion gần bằng không. Depth \code{32FC1} dùng mét; \code{16UC1} nhân 0.001. Không có fallback RGB-only đang được thực thi; ghi chú chiếu lên mặt phẳng trong YAML chỉ là hướng phát triển.

Ảnh chuyển BGR sang HSV. Hue cho đỏ là hai khoảng $0$--$10$ và $170$--$179$; vàng $20$--$38$; xanh dương $95$--$125$; xanh lá $40$--$85$; tím $130$--$165$. Mask có lọc ROI bàn và độ cao mặt vật, morphological open/close kernel 3. Contour phải có diện tích 200--1400 pixel, aspect 0.65--1.55, fill ít nhất 0.65 và kích thước vật lý phù hợp block 40 mm. Chương trình chọn candidate có score cao nhất trong các candidate đã qua lọc, không đơn thuần chọn mảng màu lớn nhất. Thiết kế giả định một block ứng với mỗi màu.
\src{src/ur3_vision_planning/worlds/vision_world.sdf.xacro}{overhead_rgbd}
\src{src/ur3_vision_planning/config/perception.yaml}{hsv, contours, depth, temporal}
\src{src/ur3_vision_planning/ur3_vision_planning/camera_perception.py}{CameraPerception.process_pending, process_pair, median_depth}
\reportfigure{images/gazebo_environment.png}{Môi trường Gazebo cần ghi nhận từ mô phỏng thực tế.}{fig:gazebo}{Chụp toàn cảnh UR3, gripper, camera, bàn, đủ năm block và ba zone; bảo đảm nhìn rõ vị trí các đối tượng.}
\reportfigure{images/camera_detection.png}{Ảnh debug nhận diện camera.}{fig:camera}{Chụp topic /vision/debug_image, thể hiện bbox hoặc nhãn đủ năm block, tọa độ và zone.}

\subsection{Tọa độ, occupancy và độ mới dữ liệu}
Tại tâm contour $(u,v)$, depth lấy median trong patch bán kính 2 pixel, tối thiểu 5 mẫu hợp lệ trong khoảng 0.05--3 m. Với nội tại CameraInfo, điểm optical được tính:
\[
\boldsymbol p_{opt}=\left(\frac{(u-c_x)d}{f_x},\frac{(v-c_y)d}{f_y},d\right)^T,
\qquad \boldsymbol p_{base}=R\boldsymbol p_{opt}+\boldsymbol t.
\]
TF được truy vấn tại timestamp RGB, target \code{base_link}. Depth nhìn mặt trên block nên chương trình giảm $z$ đi nửa chiều cao block để báo tâm vật. Tọa độ lọc bằng median lịch sử tối đa 5 mẫu. Spawn pose trong YAML không được dùng làm đầu vào vị trí nhận diện; YAML vẫn cung cấp hình học bàn, kích thước block và hình học zone.

\code{object_zone} kiểm tra toàn bộ footprint block nằm trong vùng usable $75\times75$ mm, trừ margin 1 mm, và tâm nằm trên mặt bàn với tolerance hỗ trợ 15 mm. Zone được suy ra từ tọa độ block camera so với \emph{zone đã cấu hình}; chương trình không nhận diện độc lập marker zone từ ảnh. Zone chỉ được coi empty khi camera thấy đủ năm vật; nếu thiếu vật và không có occupant nhìn thấy thì occupancy là \code{null}, status \code{unknown}. Nhiều occupant được xuất danh sách và manager từ chối.

Perception có timeout ảnh/depth/object 1 s, xuất trạng thái 5 Hz bằng steady timer để vẫn báo stale khi đồng hồ mô phỏng dừng. Vật không còn hợp lệ có \code{detected=false}, \code{position=null}, confidence 0. Confidence được tính từ fill, aspect, kích thước vật lý rồi giảm theo tuổi detection. Đây là score heuristic, không phải xác suất đã hiệu chuẩn; \code{EnvironmentPlanner.check_environment} không đặt ngưỡng confidence tối thiểu. Manager kiểm tra receive timeout 1.5 s, timestamp tiến triển, tuổi measurement 1.5 s, tuổi object 1 s và yêu cầu thấy đủ năm vật.
\src{src/ur3_vision_planning/ur3_vision_planning/camera_perception.py}{object_zone, publish_state}
\src{src/ur3_vision_planning/ur3_vision_planning/environment_manager.py}{EnvironmentPlanner.check_environment, EnvironmentManager.snapshot, on_environment}
\reportfigure{images/environment_state.png}{Bản tin camera xác định zone B bị chiếm.}{fig:state}{Chụp terminal environment_state có status ok, stale false, blue_cube.zone=zone_b và zones.zone_b.object=blue_cube; không để thông tin bí mật trong terminal.}

Ví dụ JSON đúng hình dạng dữ liệu được trình bày ở Phụ lục B. Các số trong ví dụ chỉ minh họa schema, không phải kết quả camera đo trong lần tổng hợp.

\section{LLM Planner và kiểm tra kế hoạch}
\subsection{Cấu hình 9Router và ranh giới trách nhiệm}
\code{LLMPlanner} dùng \code{urllib.request} gọi HTTP không streaming. Base URL mặc định là \code{http://localhost:20128/v1}, có thể thay bằng \code{OPENAI_BASE_URL} hoặc ROS parameter \code{openai_base_url}. Model ưu tiên ROS parameter \code{llm_model}, sau đó \code{OPENAI_MODEL}, rồi \code{LLM_MODEL}. \code{StructuredGoalPlanner} truyền default model rỗng, vì vậy thiếu model cấu hình phải bị từ chối; không tự kế thừa mặc định legacy \code{gemini/gemini-3.8-flash}.

README ghi cấu hình từng dùng là \code{oc/muse-spark-1.2-contributor-free}. Đây là bằng chứng tài liệu về cấu hình, chưa xác nhận model đang được tiến trình runtime sử dụng tại ngày tổng hợp vì không đọc cấu hình bí mật và không gọi API. POST đến \code{/chat/completions} có \code{stream=false}, \code{response_format.type=json_object}, \code{temperature=0}; timeout mặc định 45 s. GET \code{/models} chỉ có tính advisory. Credential lấy từ biến môi trường; báo cáo không trình bày giá trị credential.
\src{src/ur3_vision_planning/ur3_vision_planning/llm_planner.py}{_configured_model, LLMPlanner.__init__, _chat_completion}
\src{src/ur3_vision_planning/ur3_vision_planning/natural_language_planner.py}{StructuredGoalPlanner.__init__, structured_goal, advisory_catalog}
\todo{Xác nhận model/base URL thực sự sử dụng bằng output cấu hình đã che thông tin bí mật.}

\subsection{Schema và validator}
Schema LLM của luồng Bài 03 là đúng hai field \code{object}, \code{target_zone}. Allowlist gồm \code{red_cube}, \code{yellow_cube}, \code{blue_cube}, \code{green_cube}, \code{purple_cube} và \code{zone_a}, \code{zone_b}, \code{zone_c}. LLM không được trả tọa độ, temporary, skill, giá trị khớp hoặc trajectory; câu mơ hồ, ngoài phạm vi hoặc nhiều goal được hướng dẫn trả JSON rỗng để validator từ chối.

Kế hoạch thực thi có root đúng \code{plan} và \code{temporary_positions}. Skill được phép chỉ gồm \code{pick}, \code{place}, \code{home}. Pick cần \code{object}; place cần \code{object}, \code{destination}; home chỉ có \code{skill}. Destination là ba zone hoặc slot tạm đúng ledger do manager chọn. Tọa độ slot gồm \code{frame_id=base_link}, \code{x}, \code{y}, \code{z} hữu hạn và phải khớp chính xác selection. Kế hoạch legacy khác schema: place dùng \code{zone}, chỉ có ba màu và không giải phóng zone bị chiếm; không dùng schema đó để mô tả expanded plan Bài 03.

Validator từ chối field lạ, JSON sai/duplicate key/NaN, vật hoặc skill ngoài allowlist, pick khi đang giữ vật, place sai vật, destination bị chiếm, home khi chưa thả hoặc không ở cuối, slot không được sử dụng và chuỗi không đúng mục tiêu. Goal đã thỏa được mở rộng thành home; zone trống thành pick/place/home; zone bị chiếm thành hai cặp pick/place rồi home. Kiểm tra môi trường từ chối trạng thái partial/stale, frame sai, thiếu vật, nhiều occupant hoặc occupancy không nhất quán.

API lỗi 401/404, lỗi HTTP khác, timeout, body rỗng, content-type sai, SSE dù yêu cầu non-stream, JSON sai, thiếu message content hoặc goal sai đều dẫn đến lỗi và không sinh kế hoạch giả. Mỗi command có một POST thật, không retry tự động, không đổi model hay fallback goal. Planner mục tiêu che credential trong body preview trước khi log; vẫn cần rà soát output trước khi chụp ảnh.
\src{src/ur3_vision_planning/ur3_vision_planning/environment_manager.py}{parse_json, validate_structured_goal, DryRunPlanValidator.validate, EnvironmentPlanner.validate_plan}
\src{src/ur3_vision_planning/ur3_vision_planning/llm_planner.py}{_request_json, _message_content}
\src{src/ur3_vision_planning/ur3_vision_planning/task_validator.py}{TaskValidator}

\section{Các robot skill}
\subsection{Pick, place và home}
Trong pipeline mới, \code{pick_from_camera} kiểm tra state/object rồi truyền vị trí camera vào \code{RobotSkills.pick}. Skill mở ngón, đi đến pre-grasp, hạ Cartesian với tool hướng xuống $(x,y,z,w)=(1,0,0,0)$, đóng ngón kiểm tra tiếp xúc, tạo joint vật lý Gazebo, attach collision object trong MoveIt và nâng Cartesian. Sau nâng, pose Gazebo được so với offset tool--object lúc tiếp xúc, tolerance tối đa 15 mm theo từng trục. Đọc pose Gazebo ở đây là \emph{kiểm chứng}, không thay tọa độ camera để chọn pick target.

\code{place_destination} dùng đúng destination đã chọn. Với tâm cube cao 0.22 m, offset tool0--tâm cube 0.16 m, release clearance 5 mm và approach clearance 15 mm, final tool0 danh định 0.385 m, pre-place 0.400 m. Skill thử Cartesian approach/descent; khi planning fail có fallback pose được lập kế hoạch và kiểm tra qua adapter. Trước release, TF phải gần final pose trong 8 mm và sai số góc 0.08 rad. Gazebo detach được xác nhận trước mở ngón; sau mở, pose thực của vật được đọc để đồng bộ scene, rồi retreat. Nếu mở ngón lỗi, chương trình cố attach lại và giữ trạng thái nội bộ bảo thủ.

\code{home} của adapter lấy fresh actual joint state, lập kế hoạch đến cấu hình canonical \code{(0,-1.5707,0,0,0,0)}. Các goal tương đương chỉ được xét có giới hạn khi canonical planning failure/timeout; phải đúng FK, trong limits và collision-valid. Quỹ đạo được kiểm tra trước ExecuteTrajectory. Trường hợp đã home cũng cần state validity. Không thực hiện home khi còn object attached trong executor.

Skill hỗ trợ thực sự có gồm \code{place_at}, \code{pick_from_camera}, \code{place_destination}, \code{camera_retreat}, \code{sync_camera_scene} và các hàm motion/verification/precheck. Chúng là hàm hỗ trợ, không phải skill LLM được phép thêm vào plan. Camera retreat tại $(0.18,-0.20,0.45)$ m đưa tool khỏi ROI để quan sát lại vật.
\src{src/ur3_vision_planning/ur3_vision_planning/robot_skills.py}{pick, place_at, verify_tool_pose}
\src{src/ur3_vision_planning/ur3_vision_planning/vision_robot_skills.py}{pick_from_camera, place_destination, camera_retreat}
\src{src/ur3_vision_planning/ur3_vision_planning/continuous_motion.py}{home, move_to_pose, _continuous_cartesian_execution}

\subsection{Gripper, attach/detach và camera verification}
Gripper đang dùng là hai khớp prismatic với velocity controller. \code{command_gripper} đóng vòng qua \code{/joint_states}, đặt hai target \code{[position,0.043-position]}, vận tốc điều khiển có độ lớn 0.03 m/s. Mở dùng position 0; đóng để gắp dùng 0.043 với \code{require_object=true}. Xác nhận tiếp xúc dựa trên hai ngón dừng ở closure trung gian đối xứng và vận tốc nhỏ, không phải cảm biến lực. Vì vậy cần thực nghiệm để chứng minh gắp vật ổn định.

Plugin \code{gz::sim::systems::DetachableJoint} được nạp động sau tiếp xúc, nối \code{wrist_3_link} với \code{<object>/link}. Không nạp joint cho toàn bộ cube khi launch. \code{set_gazebo_attachment} yêu cầu acknowledgement; subscriber presence đơn thuần không chứng minh attach. Cache chỉ giữ acknowledgement đúng đối tượng/trạng thái đã quan sát trong tiến trình hiện tại và còn subscriber phù hợp. Khi transition mới timeout, cache bị invalid và dừng trước lift/release.

MoveIt biểu diễn cube bằng \code{CollisionObject} và \code{AttachedCollisionObject}; \code{update_attached_object} chuyển world/attachment và đặt touch links. Trước execution và mỗi pick, world cube được đồng bộ từ camera. Sau place và camera retreat, \code{CameraPlacementConfirmation} yêu cầu ít nhất ba measurement có timestamp khác nhau, đều sau retreat, gần destination trong 20 mm theo XY và 12 mm theo Z. Với zone phải khớp cả \code{object.zone} và \code{zone.object}; với temporary phải nằm gần điểm chọn, ngoài zone và zone nguồn trống. Timeout mặc định 20 s.
\src{src/ur3_vision_planning/config/ur3_gripper_controllers.yaml}{simple_gripper_controller}
\src{src/ur3_vision_planning/ur3_vision_planning/robot_skills.py}{command_gripper, load_gazebo_attachment_plugin, set_gazebo_attachment, update_attached_object}
\src{src/ur3_vision_planning/ur3_vision_planning/plan_execution.py}{CameraPlacementConfirmation.observe, wait_camera_confirmation}

\subsection{Đối chiếu việc thay đổi trực tiếp pose vật}
Nguồn vẫn chứa \code{set_gazebo_model_pose}, gọi service \code{/world/empty/set_pose}, và \code{restore_grasp_object} gọi hàm đó. Không thể kết luận toàn package không chứa API thay pose. Tuy nhiên \code{VisionRobotSkills.__init__} đặt \code{preserve_failed_grasp=True}; các nhánh phục hồi sau lỗi attach/lift/grasp trong \code{RobotSkills.pick} chỉ gọi \code{restore_grasp_object} khi cờ này false. Pipeline camera mới dùng adapter đó, không có lời gọi set-pose trong luồng pick/place bình thường; lỗi giữ attachment và yêu cầu phục hồi thủ công. Đây là kết luận từ đường gọi mã nguồn, chưa phải trace runtime chứng minh không có request set-pose.

Launch tạo cube bằng \code{ros_gz_sim/create} với tọa độ khởi tạo của scenario. Đặt pose khi spawn là khởi tạo môi trường, khác với teleport trong quá trình gắp/thả. Thay world collision pose và attach/detach ảo trong precheck chỉ thay \emph{MoveIt Planning Scene}, không thay vật vật lý Gazebo.
\src{src/ur3_vision_planning/ur3_vision_planning/robot_skills.py}{set_gazebo_model_pose, restore_grasp_object, pick}
\src{src/ur3_vision_planning/ur3_vision_planning/vision_robot_skills.py}{VisionRobotSkills.__init__, feasibility_precheck}
\src{src/ur3_vision_planning/launch/llm_robot.launch.py}{launch_setup}
\todo{Lưu trace runtime của luồng thực thi camera để kiểm chứng không gọi set-pose trong pick/place và ghi nhận gắp/giữ/thả thực tế.}

\section{Xử lý zone đích bị chiếm và lựa chọn vị trí tạm}
\code{EnvironmentPlanner.resolve_goal} đọc \code{state.zones[target_zone].object}. Nếu occupant khác vật đích, manager chọn vị trí tạm và thêm pick/place occupant trước pick/place mục tiêu. Với yêu cầu đưa red vào B khi blue đang ở B, chuỗi là: pick blue, place blue tại temporary đã kiểm tra, pick red, place red vào B, home. Nếu red đã ở B thì không dời lại red.

\code{temporary_positions.yaml} hiện có \textbf{12 candidate có thứ tự}; đây là tọa độ tâm cube, không phải spawn vật, không phải vị trí được LLM tạo và không bảo đảm reachability. Tầng lọc hình học yêu cầu footprint nằm trên bàn với margin 10 mm, đúng độ cao tâm 0.22 m với tolerance 5 mm, bán kính base 0.25--0.52 m, cách các cube quan sát ít nhất 80 mm, tránh marker zone có clearance và toàn bộ mặt trên còn trong camera frustum cấu hình. Không có slot đạt thì từ chối. Các kiểm tra này dùng vật camera nhận diện; không coi vật mất dấu là vị trí trống.

Tầng MoveIt xét từng slot qua IK collision-aware của pre-place và final-place khi cube được attach ảo trong scene. Slot đầu đạt cả hai được chọn; kế hoạch, tọa độ và trusted ledger được cập nhật rồi validate lại. Có thể thay slot provisional của dry-run bằng slot khác sau precheck. Temporary giữ đúng XYZ candidate, không cộng offset zone. Với zone, có thể thử offset Y nhỏ bên trong usable footprint. Do các slot có thể gần nhau, chúng là các phương án lựa chọn, không phải cam kết có thể sử dụng đồng thời.

README ghi dry-run trước đây chọn \code{temporary_1}; ghi nhận precheck ngày 03/10/2026 chọn \code{temporary_8=(0.28,-0.27,0.22)}. Không được coi temporary 8 luôn đạt trong mọi trạng thái. Ví dụ JSON ở phụ lục dùng slot này chỉ để minh họa plan sau selection, không phải chứng cứ đã dời blue.
\src{src/ur3_vision_planning/ur3_vision_planning/environment_manager.py}{candidate_rejection, free_temporary_candidates, resolve_goal, adopt_prechecked_plan}
\src{src/ur3_vision_planning/ur3_vision_planning/vision_robot_skills.py}{feasibility_precheck, destination_center}
\reportfigure{images/zone_b_occupied.png}{Trạng thái đầu của scenario zone B bị chiếm.}{fig:occupied}{Chụp Gazebo sau launch scenario zone_b_occupied, blue_cube trên zone_b; đây là bố trí spawn ban đầu.}
\reportfigure{images/expanded_plan.png}{Kế hoạch mở rộng cho red tới zone B.}{fig:plan}{Chụp câu lệnh người dùng, structured goal, occupant, FINAL EXECUTION PLAN và temporary đã qua selection; ghi rõ dry-run hay precheck-only.}
\reportfigure{images/blue_at_temporary.png}{Vật cản tại vị trí tạm: bằng chứng cần bổ sung.}{fig:temporary}{Chỉ chụp sau thực thi Gazebo: blue tại slot đã chọn và zone_b trống; kèm camera verification nếu có. Không dùng ảnh precheck ảo.}
\reportfigure{images/red_at_zone_b.png}{Vật đích tại zone B: bằng chứng cần bổ sung.}{fig:red}{Chụp red nằm trong zone_b sau chuỗi gắp/thả Gazebo; nên giữ blue ở temporary trong khung hình.}

\section{Kiểm tra tính khả thi và an toàn chuyển động}
\subsection{Precheck và giới hạn}
Motion được bật khi đồng thời \code{dry_run=false} và \code{execute_robot=true}. \code{precheck_only=true} vẫn không gửi trajectory/gripper command. \code{feasibility_precheck} xác nhận Gazebo có các model yêu cầu, TF đến tool0, dữ liệu sáu khớp, collision table và không còn attached object từ nhiệm vụ trước. Scene đồng bộ từ camera rồi lưu baseline để mô phỏng các bước tương lai.

Pre-grasp dùng IK có collision checking. Lowering và lift kiểm tra Cartesian từ seed tương ứng, \code{max_step=0.005}, fraction phải 1.0 và FK start/end phù hợp. Cube được attach ảo trước lift, detach ảo tại điểm đặt để những bước sau phản ánh zone đã giải phóng. ACM chỉ cho phép target/touch links lúc approach/lowering; reset sau attach, không bỏ collision bàn hoặc robot. Khi Cartesian fail, chẩn đoán waypoint cách nhau tối đa 2 mm lấy geometric IK rồi kiểm tra whole-robot state validity; nghiệm chẩn đoán không được dùng làm quỹ đạo hay thay fraction thiếu thành PASS.

Pre-place/final-place và camera-retreat trong precheck hiện được kiểm tra \textbf{IK tại điểm}. HOME dùng \code{GetMotionPlan} planning-only từ \textbf{seed retreat giả lập cuối}, không lấy robot thật đang home để thay start. Model/limits lấy từ URDF/SRDF runtime và MoveIt override; parameter chưa khởi tạo được phân biệt bằng type rồi đọc riêng. Description có fallback topic reliable/transient-local. Code khác SUCCESS, kể cả 99999, đều thất bại. HOME trajectory kiểm tra start, endpoint, giới hạn khớp và validity từng waypoint.

\textbf{Giới hạn quan trọng:} precheck chưa lập kế hoạch đầy đủ mọi đoạn nối đến pre-grasp, mọi approach/descent/retreat place. IK tại hai điểm không chứng minh có đường nối collision-free; Cartesian grasp/lift PASS và HOME PASS cũng không chứng minh grasp vật lý hoặc toàn nhiệm vụ thành công. Khi thực thi, từng chuyển động vẫn phải được plan và guard; nó có thể thất bại dù precheck đã PASS.
\src{src/ur3_vision_planning/ur3_vision_planning/vision_robot_skills.py}{feasibility_precheck, _ik, _cartesian_precheck, _diagnose_cartesian_waypoints, _home_model, _home_precheck}
\src{src/ur3_vision_planning/ur3_vision_planning/home_precheck.py}{home_joint_model, explicit_state, joint_values}
\src{src/ur3_vision_planning/ur3_vision_planning/vision_robot_skills.py}{_home_plan_candidate}

\subsection{Guard quỹ đạo và quản lý lỗi}
\code{ContinuousMotionMixin} giữ seed full joints và chọn winding tương đương gần seed; không dùng modulo để che jump $2\pi$. Guard kiểm tra đúng sáu tên khớp tay, dữ liệu hữu hạn, position limits, derivative vector nếu có, thời gian tăng, ít nhất hai điểm, raw first point khớp fresh actual sample trong 0.03 rad khi execution, chênh giữa hai điểm không quá 0.5 rad và endpoint FK/joint target hợp lệ. Không chỉnh sửa quỹ đạo lỗi rồi thực thi. Guard chung kiểm tra state validity ở endpoint; collision trên đường dựa vào MoveIt planner/Cartesian collision checking, không phải một chứng minh liên tục độc lập của guard. HOME có kiểm tra waypoint riêng. Chưa thấy guard kiểm tra độ lớn vận tốc/gia tốc so với dynamic limit; nó chỉ kiểm tra hữu hạn/kích thước các vector này và giới hạn vị trí.

\code{attached_object}, \code{gazebo_cube_attached}, offset grasp và acknowledgement cache theo dõi vật đang giữ. Khi bước lỗi, executor dừng cả plan, không chạy bước tiếp theo, không tự mở ngón/detach/home. Nếu motion đã bắt đầu, \code{stop_pending_motion} phát vận tốc ngón bằng 0 và cố cancel action; không xác nhận được dừng phải báo \code{STOP_NOT_CONFIRMED}. Đây là cơ chế phần mềm, không phải emergency stop phần cứng. Sau lỗi release, trạng thái nội bộ có thể bảo thủ hơn trạng thái vật lý; cần operator khôi phục simulator trước thử lại.

Trong \code{finally} của precheck, world objects, attachment và ACM được rollback, không gửi simulated joint state đến robot. \code{_verify_scene_restored} đọc lại scene kiểm tra khôi phục. Executor nhận camera mới rồi resync cả khi precheck thất bại nếu rollback đã được xác nhận. Restore/resync lỗi chặn motion; trước motion còn validate lại plan và từ chối drift camera quá 10 mm so với snapshot đã precheck.
\src{src/ur3_vision_planning/ur3_vision_planning/continuous_motion.py}{_actual_seed, _normalize_branch, _guard_trajectory, _execute_checked}
\src{src/ur3_vision_planning/ur3_vision_planning/joint_continuity.py}{branch_candidates, check_trajectory}
\src{src/ur3_vision_planning/ur3_vision_planning/vision_robot_skills.py}{stop_pending_motion, cancel_pending_precheck, _verify_scene_restored}
\src{src/ur3_vision_planning/ur3_vision_planning/plan_execution.py}{execution_enabled, execute_validated_plan}

\section{Kết quả thực nghiệm}
\subsection{Bằng chứng hiện có và mức xác nhận}
\begin{longtable}{@{}p{.25\linewidth}p{.42\linewidth}p{.26\linewidth}@{}}\toprule
Hạng mục & Bằng chứng đã đọc & Kết luận được phép\\\midrule\endhead
Build package & \path{log/build_2026-10-04_20-52-35/ur3_vision_planning/stdout.log} ghi cài script/package; stderr của lượt đó rỗng. & Có log cài đặt; không chứng minh motion.\\
Camera occupancy & README mục scenario ghi blue ở B, tọa độ camera xấp xỉ $(0.28149,0,0.22)$ m và A/C empty. & Ghi nhận tài liệu, chưa đối chiếu log camera gốc.\\
9Router dry-run & README mục natural-language ghi JSON red tới B, blue occupant, expanded plan với temporary 1; không motion. & Ghi nhận dry-run, không phải gắp/thả.\\
Precheck ba lượt & README đầu tệp ghi ngày 03/10/2026, temporary 8, Cartesian grasp/lift fraction 1.000, restore/resync PASS, HOME code 1 và 64 điểm. & Ghi nhận planning trong Gazebo; chưa có raw trace kèm theo.\\
Unit tests & Có 12 tệp test; README ghi 39/72 test tại các giai đoạn khác nhau. Chưa tìm thấy log colcon test hoặc JUnit của package trong build/log đã rà soát. & Có triển khai test; chưa xác nhận kết quả suite hiện tại.\\
Thực thi expanded plan mới & README ghi execution với trajectory guard mới chủ động chưa thử. Không tìm thấy log thành công toàn chuỗi trong các nguồn đã rà soát. & Chưa đủ bằng chứng TASK SUCCESS.\\\bottomrule
\end{longtable}

README còn ghi sai số endpoint HOME ba lượt lần lượt khoảng 0.00744, 0.00983 và 0.00955 rad, biến thiên raw branch khoảng $2.26\times10^{-14}$ rad. Các số này chỉ là \emph{số liệu được tài liệu ghi nhận}, không phải phép đo mới của báo cáo và không thay thế log gốc. Không dùng các số đó để công bố tỷ lệ thành công, thời gian chu kỳ hoặc độ chính xác đặt vật. Tài liệu có các đoạn cũ mô tả ít candidate hơn và chưa có perception; mã hiện tại phải được ưu tiên khi xác định chức năng.

Các test kiểm tra strict goal, lỗi transport, geometry temporary, thứ tự held object, rollback scene, fraction thiếu, gripper/attachment acknowledgement, HOME runtime model và continuity. Fixture service/transport không phải MoveIt solver, camera hay robot thật. Lần tổng hợp này không chạy lại pytest để giữ phạm vi đọc và phân tích mã nguồn.
\src{src/ur3_vision_planning/test/test_plan_execution.py}{test_confirmation_requires_three_distinct_post_retreat_frames}
\src{src/ur3_vision_planning/test/test_natural_language_planner.py}{test_http_error_is_not_retried_or_replaced, test_pipeline_resolves_llm_goal_with_camera_not_llm_coordinates}
\src{src/ur3_vision_planning/test/test_home_precheck.py}{test_pipeline_home_uses_last_zone_retreat_endpoint, test_home_failure_still_restores_virtual_scene}
\todo{Bổ sung log gốc dry-run, ba lượt precheck, kết quả test đúng revision và log thực thi đầy đủ. Ghi phiên bản source, model, scenario, slot chọn và chế độ chạy cùng mỗi log.}
\reportfigure{images/task_success.png}{Bằng chứng terminal thực thi toàn nhiệm vụ, chỉ bổ sung khi có log thực tế.}{fig:success}{Chỉ chèn ảnh khi các bước của cùng lượt execution Gazebo và camera verification đều đạt, cuối cùng có TASK SUCCESS. Hiện báo cáo chưa xác nhận kết quả này.}

\section{Các lỗi gặp phải và cách khắc phục}
\begin{longtable}{@{}p{.26\linewidth}p{.40\linewidth}p{.27\linewidth}@{}}\toprule
Vấn đề & Cách xử lý trong mã/tài liệu & Trạng thái bằng chứng\\\midrule\endhead
Ngón mở giao vật lân cận & Scenario occupied mới giãn bố trí, khoảng cách tâm tối thiểu 120 mm theo ghi nhận geometry; tính pre-grasp theo finger reach. & Geometry nguồn, chưa chứng minh grasp vật lý.\\
Pre-place quá cao không có IK & Approach clearance hiện 15 mm, tool pre-place danh định 0.400 m; thử candidate có giới hạn. & Mã và ghi nhận README.\\
Lowering place không đủ fraction & Release clearance 5 mm, chẩn đoán collision, không execute fraction thiếu; fallback pose vẫn qua guard. & Mã; cần log execution.\\
HOME lỗi hoặc code 99999 & Lấy start từ retreat giả lập, phân tích runtime limits, yêu cầu code SUCCESS và trajectory nối hợp lệ. & Mã, ghi nhận precheck README.\\
Typed parameter chưa khởi tạo & Đọc parameter types trước, đọc từng initialized parameter; fallback description topic. & Mã và regression fixture.\\
Attach lặp không có state mới & Dùng acknowledgement đã quan sát trong session cùng subscriber đúng kiểu; transition mới vẫn cần state. & Mã và test attachment; cần trace thực thi.\\
Jump nhánh IK/winding & Seed liên tục, equivalent branch có limits/FK/validity, guard raw trajectory delta. & Mã; README nói execution mới chưa thử.\\
Camera stale/che khuất & Invalidate detection, zone unknown, manager từ chối; retreat để thấy lại vật. & Mã, README ghi đã thử pause; cần raw log.\\
API hoặc JSON không hợp lệ & Lỗi rõ ràng, không retry/fallback giả; không motion khi goal sai. & Mã và fixture HTTP.\\\bottomrule
\end{longtable}
Đây là các vấn đề được mã và README mô tả, không phải lỗi mới quan sát trong lần tổng hợp. Không gán nguyên nhân cụ thể cho HOME chỉ từ code lỗi chung nếu thiếu contact, message/source hoặc start/goal đầy đủ.
\src{src/ur3_vision_planning/README.md}{Seed continuity; HOME feasibility; Chẩn đoán red lowering; Gazebo attachment acknowledgement}

\section{Đánh giá hệ thống}
\subsection{Các chức năng đã triển khai}
Nguồn xác nhận môi trường có UR3, gripper, RGB-D, bàn, năm cube và ba zone; perception lấy vị trí vật từ ảnh/depth; expanded planner xử lý zone occupied bằng vị trí tạm; validator kiểm tra schema và trạng thái; executor có cổng dry-run/precheck-only, attachment vật lý Gazebo, scene sync, trajectory guard và camera confirmation. Launch Bài 03 \textbf{cố định ur_type=ur3}; launch hạ tầng có hỗ trợ UR3e nhưng package chính không expose parameter chuyển loại robot. Chưa có bằng chứng kiểm thử Bài 03 với UR3e.

Luồng nên dùng để đánh giá đề bài là \code{natural_language_planner} với expanded camera pipeline. \code{skill_executor} legacy dù khởi tạo adapter motion mới vẫn dùng validator ba màu, occupancy nội bộ và pick không có camera position; không đáp ứng yêu cầu camera-first nếu chọn entrypoint đó. Chỉ dùng tên class adapter giống nhau không làm hai luồng tương đương.

\subsection{Hạn chế và rủi ro còn tồn tại}
Hệ thống hiện chỉ nhận một mục tiêu object--zone mỗi command. Zone pose là hình học cấu hình; chưa nhận diện marker độc lập. HSV phụ thuộc màu/ánh sáng, giả định một vật mỗi màu; confidence chưa có ngưỡng manager hoặc hiệu chuẩn. Yêu cầu thấy đủ năm vật giúp tránh coi vật mất dấu là chỗ trống nhưng dễ dừng khi che khuất. Candidate tạm là danh sách hữu hạn và camera frustum dùng extrinsic cấu hình; thay camera/bàn cần cập nhật và kiểm chứng TF.

Precheck chưa chứng minh toàn bộ quỹ đạo place hoặc grasp vật lý. Trong pick execution, \code{remove_world_object} gỡ target khỏi collision world trước lowering; precheck giữ target và chỉ cho touch links bằng ACM. Hai cách xử lý collision target này khác nhau và cần đối chiếu để không khẳng định execution có cùng mức collision bảo vệ như precheck. Guard không phải kiểm chứng độc lập collision liên tục hay dynamic limits. Scene publish ban đầu dùng spawn YAML, chỉ sau camera sync mới đại diện vị trí quan sát.

Nguồn vẫn còn nội dung cá nhân hóa legacy trong prompt/header của planner và executor; cần loại bỏ ở lần sửa mã được cho phép để tuân thủ việc thông tin sinh viên chỉ dùng trang bìa. Báo cáo không sử dụng nội dung đó làm yêu cầu, tính toán hay nhiệm vụ Bài 03. API set-pose phục hồi legacy vẫn tồn tại nhưng các nhánh đó bị chặn trong adapter camera hiện tại; cần kiểm tra call trace khi thực thi.

Launch chính truyền \code{world_file} vào include \code{ur_sim_moveit.launch.py}, nhưng launch trung gian trong workspace không khai báo/chuyển tiếp tường minh field này sang control. Control đọc \code{LaunchConfiguration('world_file')}; giá trị có thể đi qua launch context include. Vì vậy chưa kết luận camera world bị mất, nhưng cần xác nhận world thật được tải khi launch; không tự sửa launch hoặc coi README là chứng cứ propagation đầy đủ.

Simulator guard chỉ kiểm tra model Gazebo có mặt; không chứng minh mọi endpoint MoveIt/controller cùng thuộc simulator. Các lệnh trong báo cáo dành cho Gazebo trong ROS domain của mô phỏng, chưa có thử nghiệm phần cứng UR thật. Không dùng cơ chế guard này để công bố mức an toàn phần cứng.
\src{src/ur3_vision_planning/ur3_vision_planning/skill_executor.py}{main}
\src{src/ur3_vision_planning/ur3_vision_planning/robot_skills.py}{pick, simulator_is_ready}
\src{src/ur3_vision_planning/ur3_vision_planning/natural_language_planner.py}{goal_system_prompt, main}
\src{src/ur_simulation_gz/ur_simulation_gz/launch/ur_sim_moveit.launch.py}{launch_setup, generate_launch_description}
\todo{Xác nhận camera world được nạp, collision target khi execution, controller kết nối đúng simulator và thử nghiệm đầy đủ hai scenario.}

\section{Kết luận và hướng phát triển}
Mã nguồn đã triển khai cấu trúc LLM ở cấp ý định kết hợp camera, kế hoạch xác định, validator và MoveIt. Cơ chế giải phóng zone bằng vị trí tạm có hai tầng kiểm tra hình học và IK với vật attached, có rollback Planning Scene và kiểm tra lại camera trước motion. Thiết kế cũng tách rõ dry-run, precheck-only và execution trong Gazebo.

Bằng chứng hiện có cho phép mô tả chức năng triển khai, log build và các lượt dry-run/precheck được README ghi nhận. Chưa đủ cơ sở kết luận hệ thống đã hoàn thành toàn bộ bài thực hành hoặc thực thi thành công chuỗi giải phóng B với trajectory guard hiện tại. Hướng tiếp theo là lưu log/ảnh runtime, kiểm thử pick--hold--move--release, kiểm tra đầy đủ đường nối place trong precheck, thống nhất collision target giữa precheck/execution, loại bỏ nội dung legacy không thuộc bài và cải thiện perception khi che khuất. Mở rộng UR3e hoặc phần cứng phải là một đợt xác minh riêng.

\appendix
\section{Lệnh build và chạy}
Các lệnh dưới đây được lấy từ console scripts, launch arguments và parameter trong mã; \textbf{không được chạy trong lần tổng hợp}. Dùng terminal đã source cùng workspace/domain. Đối với credential, chuẩn bị qua môi trường terminal bằng phương thức quản lý riêng; không in credential hoặc đưa vào báo cáo. Ví dụ đặt model dưới đây dựa trên README, không xác nhận catalog hiện tại.

\subsection{Build và khởi tạo môi trường}
\begin{Code}
cd ~/ur_ws
source /opt/ros/jazzy/setup.bash
colcon build --packages-select ur3_vision_planning --symlink-install
source install/setup.bash
# Nếu package mô phỏng local chưa được build:
# colcon build --packages-select ur_simulation_gz ur3_vision_planning --symlink-install
# source install/setup.bash

# Terminal 1: chỉ chọn một scenario cho mỗi simulator
ros2 launch ur3_vision_planning llm_robot.launch.py scenario:=default
# Hoặc thay bằng:
ros2 launch ur3_vision_planning llm_robot.launch.py scenario:=zone_b_occupied
\end{Code}
Launch mặc định đã chạy perception và bridge camera, không tự chạy planner/executor. Scenario default có các zone trống; mục tiêu mẫu là một vật vào B, không phải sắp xếp cả bàn. Bố trí default còn các vật cách nhau 75 mm; cần precheck, không mặc định coi reachability giống scenario occupied đã giãn vật. Tham số GUI/RViz thực tế là \code{gazebo_gui}, \code{launch_rviz}.

\subsection{Perception và kiểm tra môi trường}
\begin{Code}
# Terminal 2, sau source ROS/workspace:
ros2 topic echo /vision/environment_state --once
ros2 run rqt_image_view rqt_image_view /vision/debug_image
# rqt_image_view là công cụ xem ảnh tùy chọn, cần có package cài sẵn.

# Chỉ dùng khi chưa có perception từ launch (tránh chạy trùng node):
ros2 run ur3_vision_planning camera_perception --ros-args -p use_sim_time:=true

# Resolve camera goal, không gọi LLM và không chuyển động:
ros2 run ur3_vision_planning environment_manager_test --ros-args \
  -p use_sim_time:=true -p object_name:=red_cube -p target_zone:=zone_b
\end{Code}

\subsection{Dry-run, precheck và execution Gazebo}
\begin{Code}
export OPENAI_BASE_URL=http://localhost:20128/v1
export OPENAI_MODEL=oc/muse-spark-1.2-contributor-free
# Credential phải đã sẵn sàng trong môi trường; không ghi giá trị ở đây.

# Dry-run LLM + camera, không motion clients:
ros2 run ur3_vision_planning natural_language_planner --ros-args \
  -p use_sim_time:=true -p dry_run:=true \
  -p command:='Đưa vật màu đỏ vào vùng B.'

# Precheck có LLM, không chuyển động:
ros2 run ur3_vision_planning natural_language_planner --ros-args \
  -p use_sim_time:=true -p dry_run:=false -p execute_robot:=true \
  -p precheck_only:=true -p command:='Đưa vật màu đỏ vào vùng B.'

# Precheck không gọi LLM; entrypoint không cho precheck_only=false:
ros2 run ur3_vision_planning motion_precheck --ros-args \
  -p use_sim_time:=true -p dry_run:=false -p execute_robot:=true \
  -p precheck_only:=true -p object_name:=red_cube -p target_zone:=zone_b

# Kiểm tra lặp planning-only nếu cần:
ros2 run ur3_vision_planning motion_precheck --ros-args \
  -p use_sim_time:=true -p dry_run:=false -p execute_robot:=true \
  -p precheck_only:=true -p object_name:=red_cube -p target_zone:=zone_b \
  -p repeat_runs:=3 -p repeat_expected_temporary:=temporary_8

# Có chuyển động trong Gazebo; chỉ dùng khi operator sẵn sàng kiểm thử:
ros2 run ur3_vision_planning natural_language_planner --ros-args \
  -p use_sim_time:=true -p dry_run:=false -p execute_robot:=true \
  -p precheck_only:=false -p command:='Đưa vật màu đỏ vào vùng B.'
\end{Code}
Precheck-only chỉ in kết quả feasibility và thông báo motion chưa bắt đầu; không phải TASK SUCCESS. Khi execution lỗi còn giữ vật, dừng executor và khôi phục mô phỏng/scenario trước retry; không dùng lại plan mới trên scene đang grasp dở. Các lệnh này không phải hướng dẫn kết nối robot phần cứng.
\src{src/ur3_vision_planning/ur3_vision_planning/natural_language_planner.py}{NaturalLanguagePlanner.__init__, main}
\src{src/ur3_vision_planning/ur3_vision_planning/motion_precheck.py}{main}

\section{JSON mẫu và danh mục minh chứng cần bổ sung}
\subsection{Structured goal và expanded plan}
Hai JSON dưới đây là mẫu theo schema, không phải response mới thu được từ API hay log thực thi:
\begin{Code}
{"object":"red_cube","target_zone":"zone_b"}

{
  "plan": [
    {"skill":"pick", "object":"blue_cube"},
    {"skill":"place", "object":"blue_cube", "destination":"temporary_8"},
    {"skill":"pick", "object":"red_cube"},
    {"skill":"place", "object":"red_cube", "destination":"zone_b"},
    {"skill":"home"}
  ],
  "temporary_positions": {
    "temporary_8":{"frame_id":"base_link","x":0.28,"y":-0.27,"z":0.22}
  }
}
\end{Code}
Expanded plan chỉ hợp lệ nếu temporary 8 là selection tin cậy của manager/precheck trong đúng snapshot, không được chép mẫu để bỏ qua kiểm tra.

\subsection{Environment state}
Ví dụ đầy đủ năm object/ba zone sau đây minh họa schema. Timestamp và confidence là giá trị giả định để giải thích field, không dùng làm số liệu thực nghiệm; pixel/bbox trong bản tin thực có thể là object khi detection hợp lệ.
\begin{Code}
{
 "stamp":100.0,"published_at":100.1,"frame_id":"base_link",
 "stale":false,"status":"ok","reason":"ok",
 "objects":{
  "red_cube":{"detected":true,"stale":false,
   "position":{"x":0.38,"y":-0.15,"z":0.22},"zone":null,
   "confidence":0.9,"age_s":0.1,"stamp":100.0,
   "pixel":{"u":300,"v":240},"bbox":{"x":289,"y":229,"width":22,"height":22}},
  "yellow_cube":{"detected":true,"stale":false,
   "position":{"x":0.40,"y":-0.03,"z":0.22},"zone":null,
   "confidence":0.9,"age_s":0.1,"stamp":100.0,"pixel":null,"bbox":null},
  "blue_cube":{"detected":true,"stale":false,
   "position":{"x":0.28,"y":0.0,"z":0.22},"zone":"zone_b",
   "confidence":0.9,"age_s":0.1,"stamp":100.0,"pixel":null,"bbox":null},
  "green_cube":{"detected":true,"stale":false,
   "position":{"x":0.35,"y":0.21,"z":0.22},"zone":null,
   "confidence":0.9,"age_s":0.1,"stamp":100.0,"pixel":null,"bbox":null},
  "purple_cube":{"detected":true,"stale":false,
   "position":{"x":0.38,"y":0.09,"z":0.22},"zone":null,
   "confidence":0.9,"age_s":0.1,"stamp":100.0,"pixel":null,"bbox":null}
 },
 "zones":{
  "zone_a":{"occupied":false,"object":null,"objects":[],
   "stale":false,"status":"empty"},
  "zone_b":{"occupied":true,"object":"blue_cube","objects":["blue_cube"],
   "stale":false,"status":"occupied"},
  "zone_c":{"occupied":false,"object":null,"objects":[],
   "stale":false,"status":"empty"}
 }
}
\end{Code}

\subsection{Ảnh và xác nhận còn thiếu}
Ảnh được đặt trong \path{reports/images/} khi biên dịch từ \path{reports/}. Mỗi figure dùng \code{\IfFileExists}; nếu thiếu ảnh, khung ghi rõ nội dung cần chụp vẫn được biên dịch.
\begin{longtable}{@{}p{.49\linewidth}p{.44\linewidth}@{}}\toprule
Tên ảnh & Phần sử dụng\\\midrule\endhead
\path{images/gazebo_environment.png} & Camera: toàn cảnh môi trường.\\
\path{images/camera_detection.png} & Camera: debug nhận diện năm block.\\
\path{images/zone_b_occupied.png} & Zone occupied: bố trí ban đầu.\\
\path{images/environment_state.png} & Camera: JSON occupancy từ terminal.\\
\path{images/expanded_plan.png} & Zone occupied: command và plan đã validate.\\
\path{images/blue_at_temporary.png} & Zone occupied: kết quả dời blue thực tế.\\
\path{images/red_at_zone_b.png} & Zone occupied: kết quả đặt red thực tế.\\
\path{images/task_success.png} & Kết quả thực nghiệm: chỉ khi có log execution thật.\\\bottomrule
\end{longtable}
\todo{Bổ sung raw log và ảnh của cùng một lượt thực thi; xác nhận slot thực được chọn, model thực dùng, revision source, camera verification, attach/detach và kết quả từng step.}
\todo{Chạy và lưu test suite hiện tại, kiểm tra scenario default và occupied, xác nhận không có set-pose trong thao tác và không dùng vị trí spawn làm target pick.}
\todo{Môi trường tổng hợp không tìm thấy executable XeLaTeX, nên chưa kiểm tra biên dịch. Khi có XeLaTeX, chạy hai lần từ reports/: xelatex -interaction=nonstopmode -halt-on-error main.tex.}
\end{document}
